\chapter{العصبون ليس جهازًا واحدًا}
% full version: chapters/18_eetypes.tex

في الفصل الأول فرزنا العصبونات بحسب المادة التي تقذفها. أما المهندس فكان سيفرزها على نحو آخر: أي جهاز يعمل كل منها. ويتبين أن تحت المظهر الواحد يختبئ صندوق كامل من القطع الإلكترونية.

\section*{صندوق القطع الإلكترونية}

الدلو المثقوب من الفصل الثاني مجمِّع تسريبي: يجمع الإشارات خلال نافذة قصيرة. وحين تصغر نافذته جدًا يتحول إلى كاشف تزامن. والعصبون الذي لا يستجيب إلا للتغيرات ثم يصمت فورًا مرشح يمرر التغيرات. والعصبون الذي يعمل من تلقاء نفسه ميقاتية، أي مذبذب؛ وإذا كان المدخل يعدّل إيقاعه، فهو مذبذب يتحكم فيه الجهد. والعصبون الذي يطلق رشقات من النقرات مولّد رشقات يكتب ”الشرطة“. والعصبون الذي ينقر بإيقاع الموجة الصوتية كاشف طور.

\pencilfig{18}{\pencilart{18}}{تحت المظهر الواحد يختبئ صندوق كامل من القطع الإلكترونية: مقاومة، ومكثف، وملف، ومذبذب.}

وثمة أجهزة لم نتحدث عنها من قبل. \emph{المرنان} يشبه الأرجوحة: يستجيب أفضل ما يستجيب لدفعات بإيقاع معين، كالأرجوحة التي لا تتأرجح إلا إذا دفعتها في الوقت المناسب. وداخل العصبون يؤدي دور النابض تيار بطيء يقاوم التغيرات. و\emph{المجمِّع بلا تسريب} يحفظ القيمة طويلًا، كمكثف لا يتسرب منه شيء: هكذا تعمل الشبكة التي تثبّت موضع العينين حين تنظر جانبًا. لكن لذلك يجب أن تُضبط التغذية الراجعة ضبطًا شبه مثالي، وهو ضبط مكلف. و\emph{المزلاج} عصبون تنقله دفعة قصيرة إلى عمل طويل إلى أن يُطفأ، كمفتاح ذي تثبيت؛ وعند العصبونات التي تتحكم في العضلات يشغّل السيروتونين هذا النمط.

\section*{جهاز حرباء}

أهم اكتشاف في هذا الفرز: هذه ليست قطعًا مختلفة، بل \emph{أنماط} مختلفة للقطعة نفسها. العصبون نفسه قد يكون مجمِّعًا في النهار ومولّد رشقات في الليل، بحسب ما تقوله له المحطات الإذاعية. سيسمي المهندس الدماغ مصفوفة منطقية قابلة لإعادة البرمجة: الدائرة واحدة، ووظائف القطع تغيّرها الإعدادات.

أن العصبونات قادرة على العمل بكل هذه الأنماط أمر مقيس. أما كيف يستخدم الدماغ تبديل الأنماط في الحساب ففي معظمه فرضية.

\fullversion{18}
